\subsection{Lessons from Human Child Development and Exploration}
\label{sec:4.3.2}
Developmental psychology has one result that transfers cleanly enough to be worth building on. Vygotsky's zone of proximal development holds that learning happens fastest at the boundary of what a learner can already do, in company with someone slightly ahead of them \autocite{vygotsky1978mind}. The design consequence is a curriculum rather than a benchmark: environments pitched just past current capability, and other agents — human or artificial — whose competence is a step ahead.

Three cautions belong with it, because the analogy is doing more work here than it can bear. A child's exploration is bounded by a body, by exhaustion, and by adults who intervene, and none of those bounds exists for a system that can run a million episodes overnight; the boundary has to be built. The developmental sequence is not a law, for the reasons section 5.1.1's box sets out. And autonomy inside limits is the whole difficulty rather than a design principle: how much room a system gets to explore is the same question chapter 3 asks about how much room it gets to refuse, arriving here in a friendlier vocabulary.
